\subsection{Multi-Wavelength Pulse Processing}
For each device, visible and infrared pulse rates are estimated around the scheduled reference time using a window of up to 12~s, with at least 7.5~s and 150 samples required. Timestamps are converted to seconds, sorted, and deduplicated; signals are interpolated to a uniform grid using the median interval, with inferred sampling rate clipped to 20--500~Hz. A third-order 0.55--4~Hz Butterworth filter is applied forward and backward after detrending. Welch spectra use up to 8~s segments (at least 256 samples when available), with 50\% overlap. The dominant frequency is sought in 0.65--3.5~Hz; spectral concentration $c$ is the fraction of pulse-band power within 0.18~Hz of that peak.

An autocorrelation pulse estimate searches lags corresponding to 39--210~BPM. Let $a$ be the selected autocorrelation peak clipped to $[0,1]$, and $h_F,h_A$ the spectral and autocorrelation rates. When $|h_F-h_A|\leq15$~BPM, the channel estimate is $(c h_F+a h_A)/(c+a)$; otherwise the estimate with the larger reliability component is used. Channel quality is
\begin{align}
q&=\operatorname{clip}\!\left(\sqrt{ca}\,g,10^{-4},1\right),\\
g&=\begin{cases}1,&|h_F-h_A|\leq15,\\
e^{-|h_F-h_A|/30},&\text{otherwise.}\end{cases}
\end{align}
The tone-blind fusion estimate is $(q_Vh_V+q_Ih_I)/(q_V+q_I)$. The optional tone-dependent arm replaces the weights by $q_Ve^{-\alpha z}$ and $q_Ie^{\alpha z}$, with $z=\operatorname{clip}[(\mathrm{MST}-1)/9,0,1]$. Candidate $\alpha$ values are $\{0,\pm0.25,\pm0.5,\pm0.75,\pm1,\pm1.5,\pm2,\pm3\}$. Four participant folds within each outer training set score these fixed candidate rules. The selected nonzero exponent must improve mean fold MAE by at least 0.10~BPM and improve every inner fold; otherwise $\alpha=0$. The seed is 20260827, with repeat/fold offsets defined in the companion script. These methods estimate pulse rate, not BP.

\subsection{VitalDB Extraction and Model Selection}
Eligible cases have all three required tracks and a case duration of at least 1,800~s. Cases are deterministically shuffled with seed 20260827; positions 1--60 define development and 61--300 define the candidate holdout. The saved audit contains 55 development cases/subjects and 228 eligible holdout cases from 227 subjects, with no subject overlap. Eligibility and sampling are retrospective, and selecting windows across the completed record is not an online scheduling policy.

Candidate centers start no earlier than 120~s, proceed in 30-s increments, and require at least 900~s of usable overlapping record. PPG is retained at approximately 100~Hz. A 20-s window must contain at least 75\% of the expected samples and 85\% finite values; gaps are interpolated. Values are clipped to the median plus/minus eight robust SD estimates, detrended, and filtered with a third-order 0.5--8~Hz bandpass (upper cutoff limited by sampling rate). Pulse quality is spectral concentration multiplied by $\exp[-\min(\mathrm{interval\ CV},2)]$; windows below 0.04 are excluded. Pulse frequency search, spectral concentration, and beat-interval limits follow the companion extraction code. This quality filter is held fixed, but a no-quality-filter ablation is not reported.

Reference BP is the median of at least three plausible values within $\pm5$~s of the center. SBP must lie in 70--220~mmHg and DBP in 35--130~mmHg; local robust SD must not exceed 8 and 6~mmHg, respectively. Pulse pressure must lie in 20--130~mmHg. These reference-based exclusions do not use prediction errors but restrict the evaluated BP distribution. Cases require at least eight valid windows. The first three define median BP and feature anchors; up to 24 later windows are selected using evenly spaced indices over the valid future windows.

The context model uses anchor BP, $\log(1+\mathrm{elapsed\ seconds})$, age, sex, and preoperative hypertension (missing hypertension is coded as zero). The PPG model adds current and calibration-relative optical features listed by the \texttt{columns} function in the companion script. Candidate estimators comprise standardized ridge regression with 13 logarithmically spaced penalties from $10^{-3}$ to $10^3$ and three absolute-loss gradient-boosting configurations. Development-only five-fold participant validation selects estimators, followed by a fixed blend weight from $\{0,0.25,0.5,0.75,1\}$; a nonzero weight requires at least 0.10~mmHg improvement over context. Development selection is not independently nested across the estimator and blend choices, but holdout observations are excluded from both. Delta predictions are clipped at training-target 1st/99th percentiles.

Both endpoints selected ridge context models with penalty approximately 31.6228. The SBP PPG model uses absolute-loss boosting, learning rate 0.03, 15 leaves, regularization 5, 250 iterations, and minimum leaf size 20; the DBP PPG model uses ridge with penalty approximately 31.6228. Both selected a constant 0.5 PPG weight. The saved predictions satisfy this fixed-blend identity to numerical precision. Errors are averaged over all windows from all cases of a subject before subject bootstrap; the subject contributing two cases remains one bootstrap unit.

\subsection{Audit Scope and Reproduction Status}
The independent-dataset saved-result audit uses 10,000 participant bootstrap resamples with seed 20260910 for the pulse comparisons and verifies window uniqueness, case/subject counts, development/holdout separation, the fixed VitalDB blend, and the pulse-fusion equation. The AuroraBP analysis was separately rerun: the context arms and all three calibration doses were refitted across three repeats of five participant folds. The package records hashed feature inputs, software versions, selected features and gates, row-level outer-fold predictions, and paired participant contrasts. The historical subgroup model was also refitted. The rerun starts from feature banks; raw-waveform extraction is outside its scope.
